\chapter{Trascrizione interviste con i soci aziendali}\label{app:risultati-interviste-proprieta} % Verificare se "proprietà" è il termine voluto
In questa appendice sono riportate le trascrizioni delle interviste effettuate ai soci dell'azienda,
al fine di raccogliere informazioni sulle esigenze e le problematiche attuali del sistema di preventivazione.
\section{Situazione attuale e problematiche}
Attualmente l'ufficio commerciale utilizza un database Microsoft Access, in uso da
oltre 20 anni.
\begin{itemize}
      \item \textbf{Punti di forza:} Estrema velocità di inserimento per utenti esperti; personalizzazione
            storica delle stampe e delle logiche.
      \item \textbf{Volume dati:} Circa 6 offerte/giorno per utente (commesse) e 18/giorno (ricambi/\ac{RA}).
            Totale annuo stimato: \(>2200\) offerte commesse.
      \item \textbf{Limitazioni critiche:}
            \begin{itemize}
                  \item Mancanza di integrazione con il gestionale \ac{ERP} (Ad Hoc).
                  \item Necessità di reinserimento dati manuale da Offerta a Ordine (doppio lavoro).
                  \item Gestione prezzi manuale o basata sull'esperienza (non collegata al costo
                        industriale aggiornato).
            \end{itemize}
\end{itemize}

\section{Criticità rilevate}
\subsection{Disallineamento tecnico-commerciale}
Attualmente, una volta firmato il contratto commerciale, le modifiche tecniche (\ac{PM}/\ac{UT})
non vengono tracciate nel sistema di preventivazione.
\begin{itemize}
      \item \textbf{Conseguenza:} Per replicare un impianto o vendere ricambi corretti, il commerciale
            deve consultare gli schemi CAD/Tecnici e il gestionale,
            poiché l'offerta originale in Access diventa obsoleta/inaccurata dopo la conferma d'ordine.
      \item \textbf{Rischio:} Vendita di ricambi errati o incoerenze nelle offerte copiando vecchie offerte
            per recuperare informazioni, condizioni e note.
\end{itemize}

\subsection{Inefficienza Gestione Ricambi (Spare Parts)}
Il flusso per i ricambi è eccessivamente complesso rispetto al valore economico.
\begin{itemize}
      \item \textbf{Situazione attuale:} L'Area Manager deve aprire una nuova offerta da zero o copiare
            vecchie offerte, richiedere
            quotazioni interne se non già presenti, attendere risposte e infine emettere l'offerta.
            Dopo la conferma del cliente, deve trasformarla in ordine.
      \item \textbf{Collo di bottiglia:} Troppi passaggi (Commerciali \(\to\) Ufficio Tecnico \(\to\) Pricing Office \(\to\) Offerta).
      \item \textbf{Obiettivo:} Ridurre drasticamente i tempi per rendere profittevole la vendita proattiva,
            che ad oggi non viene effettuata.
\end{itemize}

\section{Requisiti funzionali emersi}

\subsection{Gestione offerte e configuratore}
Il sistema deve guidare l'utente, specialmente se junior, nella composizione
dell'impianto.
\begin{itemize}
      \item \textbf{Selezione guidata:} Selezionando un macro-componente (es. ``Trasporto Pneumatico'' o ``STX''),
            il sistema deve proporre in automatico gli accessori correlati (filtri, valvole, compressori) tramite menu a tendina.
      \item \textbf{Alert di completezza:} Avvisi automatici in caso di macchine essenziali mancanti
            (es. ``Hai inserito il trasporto, ma manca il compressore'').
      \item \textbf{Template per applicazione:} Possibilità di richiamare configurazioni standard (Template di impianto)
            basate sul prodotto finale del cliente (es. ``Linea Pan di Spagna'', ``Linea Biscotti'') e non solo sul nome cliente.
\end{itemize}

\subsection{Listini e pricing}
\begin{itemize}
      \item \textbf{Automazione costi:} Il prezzo di vendita deve scaturire automaticamente dal costo industriale
            aggiornato (integrazione con il gestionale), applicando regole di ricarico predefinite.
      \item \textbf{Eliminazione colli di bottiglia:} Il commerciale non deve dover chiedere quotazioni interne per
            componentistica standard.
\end{itemize}

\subsection{Gestione flusso obbligatorio ordine e revisioni}
\begin{itemize}
      \item \textbf{Integrazione ERP:} Passaggio automatico da ``Offerta Confermata''
            a ``Ordine Gestionale'' senza reinserimento dati.
      \item \textbf{Doppio contratto (Versionamento):}
            \begin{enumerate}
                  \item \textbf{Contratto commerciale:} La versione firmata dal cliente (congelata).
                  \item \textbf{Contratto tecnico (``As-Built''):} Una copia viva che viene modificata dai
                        PM/Ufficio Tecnico durante la realizzazione.
                        Questa diventerà la base per le future assistenze e i ricambi.
            \end{enumerate}
\end{itemize}

\subsection{Distinzione flussi di lavoro}
Il software deve distinguere a monte la tipologia di richiesta per instradare
il lavoro:
\begin{enumerate}
      \item \textbf{Nuovi impianti / commesse:} Flusso standard completo.
      \item \textbf{Ampliamenti:} Collegati a commesse esistenti ma gestiti dai commerciali.
      \item \textbf{Ricambi (spare parts):} Flusso semplificato e rapido. Accesso allo storico installato.
\end{enumerate}

\section{Requisiti non funzionali}
\begin{itemize}
      \item \textbf{Modalità offline:} Fondamentale per gli Area Manager in viaggio in zone
            con scarsa connettività.
            Sincronizzazione al rientro o via \ac{VPN}.
      \item \textbf{Multilingua:} Gestione avanzata delle descrizioni (necessità di una via di mezzo tra la descrizione
            ``estesa'' attuale, troppo lunga, e quella ``breve'', troppo criptica).
      \item \textbf{Velocità:} L'interfaccia non deve rallentare l'operatività dell'utente.
\end{itemize}